\documentclass[11pt]{article}

\usepackage{amsmath,amssymb}
\usepackage{xurl}
\usepackage[colorlinks=true,linkcolor=blue,citecolor=blue,urlcolor=blue]{hyperref}

% Shared commands for notes in docs/paper-gaps/.

\DeclareUnicodeCharacter{2080}{\ifmmode{}_{0}\else\textsubscript{0}\fi}
\DeclareUnicodeCharacter{2081}{\ifmmode{}_{1}\else\textsubscript{1}\fi}
\DeclareUnicodeCharacter{2082}{\ifmmode{}_{2}\else\textsubscript{2}\fi}
\DeclareUnicodeCharacter{2099}{\ifmmode{}_{n}\else\textsubscript{n}\fi}

\newcommand{\Lean}{\textsc{Lean}}
\ExplSyntaxOn
\NewDocumentCommand{\leanid}{m}{
  \ifmmode
    \textsf{\detokenize{#1}}
  \else
    \group_begin:
    \urlstyle{sf}
    \tl_set:Nn \l_tmpa_tl {#1}
    \tl_replace_all:Nnn \l_tmpa_tl {\_} {_}
    \exp_args:NV \path \l_tmpa_tl
    \group_end:
  \fi
}
\ExplSyntaxOff
\providecommand{\tr}{\operatorname{tr}}
\newcommand{\tnleanIssueBase}{https://github.com/LionSR/TNLean/issues/}
\newcommand{\tnleanPrBase}{https://github.com/LionSR/TNLean/pull/}
\newcommand{\tnleanDocBase}{https://sirui-lu.com/TNLean/docs/}
\newcommand{\ghissue}[1]{\href{\tnleanIssueBase#1}{GitHub issue \##1}}
\newcommand{\ghpr}[1]{\href{\tnleanPrBase#1}{PR \##1}}
\newcommand{\doclink}[1]{\href{\tnleanDocBase#1.html}{\path{#1}}}

% Dirac notation and the identity operator, as in blueprint/src/macros/common.tex.
% \providecommand keeps note-local definitions authoritative.
\usepackage{dsfont}
\providecommand{\ket}[1]{|#1\rangle}
\providecommand{\bra}[1]{\langle #1|}
\providecommand{\braket}[2]{\langle #1 | #2 \rangle}
\providecommand{\ketbra}[2]{|#1\rangle\!\langle #2|}
\providecommand{\Id}{\mathds{1}}
\providecommand{\one}{\mathds{1}}
\providecommand{\C}{\mathbb{C}}
\providecommand{\MN}[1]{M_{#1}(\C)}
\providecommand{\MD}{M_D(\C)}

% External referencing for paper sources.  The build helper writes sanitized
% auxiliary files under build/paper-aux/ and excludes duplicated source labels.
\usepackage{xr}

% Import a generated paper auxiliary file with a caller-supplied label prefix.
% The candidate paths support builds from docs/paper-gaps/, the repository root,
% and build/paper-aux/ itself.
\newcommand{\importpaperaux}[2]{%
  \IfFileExists{../../build/paper-aux/#2.aux}{%
    \externaldocument[#1]{../../build/paper-aux/#2}%
  }{%
    \IfFileExists{build/paper-aux/#2.aux}{%
      \externaldocument[#1]{build/paper-aux/#2}%
    }{%
      \IfFileExists{#2.aux}{%
        \externaldocument[#1]{#2}%
      }{}%
    }%
  }%
}

% General external reference macros
\newcommand{\paperref}[2]{\ref{#1-#2}}
\newcommand{\papereqref}[2]{(\ref{#1-#2})}

% CPSV16 source (1606.00608 / MPDO-22-12-17-2.tex) external document and shortcuts
\importpaperaux{cpsv16-}{1606.00608/MPDO-22-12-17-2-tnlean}
\newcommand{\cpsvref}[1]{\paperref{cpsv16}{#1}}
\newcommand{\cpsveqref}[1]{\papereqref{cpsv16}{#1}}

% Verdict marker: \gapnote{<kind>}{<status>}. Kind is the policy
% classification (clarification, local-correction, scope-restriction,
% unfaithful, false-source, open-gap); status is open, wip (elimination
% underway), resolved, or historical. Machine-read by the site index;
% prints a verdict line.
\newcommand{\gapnote}[2]{\par\noindent\textbf{Verdict:} #1 (#2).\par}


\title{Corner-Support Restriction in the Weighted Fixed-Point Star-Algebra}
\author{}
\date{\today}

\begin{document}
\maketitle

\section{The assertion}

Corollary 6.7 of \emph{Quantum Channels \& Operations} (Wolf 2012) states: for a
trace-preserving, positive, linear map $T$ on $M_d(\mathbb{C})$ whose adjoint satisfies
the Schwarz inequality, and for any density matrix $\rho$ which is a \emph{maximum-rank}
fixed point of $T$, the set
\[
    \rho^{-1/2}\,\mathcal{F}_T\,\rho^{-1/2},
    \qquad \mathcal{F}_T = \{X \mid T(X) = X\},
\]
is a $*$-algebra, with the inverse square root taken on the support of $\rho$.

The corollary rests on the maximal-support property of fixed points stated earlier in
Section~6.4 of the reference: every fixed point of $T$ has support and range within the
support of the maximum-rank fixed point, so conjugation by $\rho^{-1/2}$ loses nothing.

\section{The formalization}

Two layers are established, both for a Kraus map (the completely positive case of the
Schwarz hypothesis, as in the other results of this section).

For an \emph{arbitrary} positive semidefinite fixed point $\rho$ with support
projection $Q$ (\leanid{Kraus.weightedCornerFixedPointsStarSubalgebra} and
\leanid{Kraus.exists\_weightedCorner\_sqrt\_eq\_of\_fixedPoint} in
\leanid{TNLean/Channel/FixedPoint/WeightedCornerFixedPoints.lean}), the set
\[
    \rho^{-1/2}\,\{X \mid T(X) = X,\ Q X Q = X\}\,\rho^{-1/2}
\]
is a $*$-subalgebra of the corner algebra $Q M_D(\mathbb{C}) Q$; for a non-maximal
$\rho$ the conjugated set is restricted to the fixed points supported on the support
of $\rho$. The positive definite case without any support restriction is
\leanid{Kraus.weightedFixedPointsStarSubalgebra}.

At a fixed point of maximal support the restriction disappears
(\leanid{Kraus.exists\_maximalSupport\_weightedCorner\_sqrt\_eq} in
\leanid{TNLean/Channel/FixedPoint/MaximalSupport.lean}): there is a positive
semidefinite fixed point $\rho_0$ at which the conjugated set is
\[
    \rho_0^{-1/2}\,\{X \mid T(X) = X\}\,\rho_0^{-1/2},
\]
the set of the corollary, with the inverse square root taken on the support of
$\rho_0$.

\section{The maximal-support property}

The removal of the restriction rests on the maximal-support property
(\leanid{Kraus.exists\_maximalSupport\_fixedPoint}): there is a positive semidefinite
fixed point $\rho_0$ whose support projection $Q_0$ satisfies $Q_0 X Q_0 = X$ for every
fixed point $X$. The witness is the sum of the positive parts of the Hermitian and
anti-Hermitian components of a basis $X_0, \ldots, X_{m-1}$ of the fixed-point space
(the reference instead takes the image of the identity under the projection onto the
fixed-point space). Writing
$H_{1,j} = X_j + X_j^{\dagger}$ and $H_{2,j} = i(X_j - X_j^{\dagger})$ for the fixed
Hermitian components and $H_{k,j} = P^{+}_{k,j} - P^{-}_{k,j}$ for their positive parts
(\leanid{IsChannel.posSemidef\_parts\_of\_hermitian\_fixedPoint}), the witness
\begin{align*}
    \rho_0 &= \sum_{j=0}^{m-1}
        \bigl(P^{+}_{1,j} + P^{-}_{1,j} + P^{+}_{2,j} + P^{-}_{2,j}\bigr)
\end{align*}
dominates every summand in the Loewner order, domination transfers the support
(\leanid{Kraus.stationaryProj\_absorb\_of\_le}, from
$0 \preceq (\mathbf 1-Q_0)P(\mathbf 1-Q_0) \preceq (\mathbf 1-Q_0)\rho_0(\mathbf 1-Q_0) = 0$),
and linearity extends the absorption from the parts to every fixed point.

\section{Verdict}

The restriction is eliminated. The maximal-support property holds at the constructed
witness (\leanid{Kraus.exists\_maximalSupport\_fixedPoint}), and the transfer to every
maximum-rank fixed point is
\leanid{Kraus.maximalSupport\_of\_maximalRank} in
\leanid{TNLean/Channel/FixedPoint/MaximalRank.lean}: for any positive semidefinite fixed
point $\rho$ whose rank bounds the rank of every fixed-point density matrix, the support
projection of $\rho$ carries every fixed point, because the support projection of a
positive semidefinite matrix has the rank of the matrix and its trace is that rank, so
the two support projections have equal traces, one absorbs the other, and they coincide.
The conjugated set $\rho^{-1/2}\,\{X \mid T(X) = X\}\,\rho^{-1/2}$ of the corollary is
realized at every such $\rho$
(\leanid{Kraus.exists\_weightedCorner\_sqrt\_eq\_of\_maximalRank}), with the inverse
square root taken on the support of $\rho$. Unit trace of $\rho$ itself is not needed for
the conclusion and is not assumed; the maximality hypothesis quantifies over fixed-point
density matrices, as in the source. The remaining scope convention is the Kraus
(completely positive) case of the source's Schwarz hypothesis, shared by the other
formalized results of the section.

\end{document}
